La datació per carboni 14 és una eina molt útil per a estimar l'edat de restes òssies o fòssils. Aquesta tècnica es basa en el cicle següent:
\begin{itemize}[label={---},itemsep=0pt]
  \item Un nucli de nitrogen 14, $^{14}_{7}\mathrm{N}$, captura un neutró provinent de rajos còsmics de l'espai, allibera un protó i es converteix en carboni 14.
  \item El carboni 14 forma molècules de $\mathrm{CO_2}$ que són absorbides per les plantes.
  \item Quan els animals mengen, les plantes incorporen el carboni 14.
  \item Quan un animal mor, ja no incorpora més carboni 14. A partir d'aquest punt el contingut de carboni 14 disminueix progressivament i es converteix en nitrogen 14.
\end{itemize}

\begin{apartat}{1,25}
Escriviu la reacció nuclear mitjançant la qual el nitrogen 14 es transforma en carboni 14 per l'efecte dels rajos còsmics. Justifiqueu si la reacció absorbeix o allibera energia.
\end{apartat}

\begin{apartat}{1,25}
Al laboratori es comparen dues mostres òssies d'elefant. La mostra $A$ és d'un individu mort recentment i la mostra $B$ té datació desconeguda. Sabent que la mostra $B$ conté un 23\,\% menys de carboni 14 que la mostra $A$, quina edat té aquesta mostra?
\end{apartat}

\dades{%
  Període de semidesintegració del carboni 14: 5\,730 anys.\\
  Masses (en unitats de massa atòmiques):\\[4pt]
  \begin{tabular}{|c|c|c|c|}
    \hline
    \textit{Neutró} & \textit{Protó} & \textit{Nitrogen 14} & \textit{Carboni 14}\\
    \hline
    1,008\,664\,9 u & 1,007\,276\,47 u & 14,003\,074 u & 14,003\,241 u\\
    \hline
  \end{tabular}}
